\chapter{Phenylketonuria (PKU)}

\section*{What Is This Test?}
Phenylketonuria (PKU) screening is a newborn screening test that measures blood phenylalanine (Phe) levels to detect phenylketonuria, an autosomal recessive inborn error of metabolism caused by a deficiency of the enzyme phenylalanine hydroxylase (PAH). PAH normally converts the essential amino acid phenylalanine to tyrosine. In PKU, phenylalanine accumulates in the blood and brain, causing progressive, irreversible intellectual disability (severe intellectual disability if untreated, IQ <50), microcephaly, seizures, autism-like behavior, musty/mousy body odor (from phenylacetic acid in sweat and urine), hypopigmentation (fair skin, light hair, blue eyes due to tyrosine deficiency reducing melanin synthesis), and eczema \cite{scriver2001metabolic}. Newborn screening is performed on dried blood spots collected by heel stick onto filter paper (Guthrie card) at 24--48 hours of life (after sufficient protein feeding to detect elevated phenylalanine). PKU is included in the Recommended Uniform Screening Panel (RUSP) and is screened in all U.S. states and most developed countries \cite{acmg2012newborn}. Detection in the newborn period allows immediate initiation of a phenylalanine-restricted diet and medical formula (special phenylalanine-free amino acid formula), which prevents intellectual disability and allows normal neurodevelopment. PKU has an incidence of approximately 1:10,000 to 1:15,000 live births in the United States (higher in Caucasians and Native Americans, lower in African Americans and Japanese).

\section*{Alternative Names}
\begin{itemize}
  \item PKU Screening
  \item Newborn Screening for PKU
  \item Phenylalanine Screening
  \item Guthrie Test (historical)
  \item Blood Phenylalanine Level
  \item PAH Deficiency
  \item Hyperphenylalaninemia Screen
\end{itemize}
\section*{Principle}
Newborn screening for PKU is performed on dried blood spots (DBS) collected on filter paper. Phenylalanine is measured by tandem mass spectrometry (MS/MS), which simultaneously quantifies phenylalanine, tyrosine, and the phenylalanine/tyrosine (Phe/Tyr) ratio. The Phe level is elevated in PKU; the Phe/Tyr ratio increases the specificity of screening. Confirmatory testing in a symptomatic infant or following a positive newborn screen uses plasma quantitative amino acid analysis by ion-exchange chromatography or HPLC with post-column ninhydrin derivatization and spectrophotometric detection. A plasma phenylalanine >120 mcmol/L (2 mg/dL) is abnormal; classical PKU typically presents with a plasma Phe >1,200 mcmol/L (20 mg/dL). The original method was the Guthrie bacterial inhibition assay introduced in 1961: dried blood spots on filter paper were placed on an agar plate with Bacillus subtilis spores containing a growth inhibitor (beta-2-thienylalanine); phenylalanine from the blood overcame the inhibition, and bacterial growth around the spot was proportional to phenylalanine concentration. This method has been largely replaced by MS/MS. Molecular genetic testing of the PAH gene (chromosome 12q23.2, >1,000 known mutations) can confirm the diagnosis and predict the severity (dietary Phe tolerance).

\section*{History}
PKU was first described in 1934 by the Norwegian physician Asbjorn Folling, who detected phenylpyruvic acid (a phenylalanine metabolite) in the urine of two siblings with severe intellectual disability. The autosomal recessive inheritance was established shortly thereafter. In 1953, the German pediatrician Horst Bickel developed the first dietary treatment for PKU (a low-phenylalanine formula). In 1961, Robert Guthrie developed the bacterial inhibition assay (Guthrie test) for mass newborn screening \cite{guthrie1963simple}. Massachusetts was the first state to mandate newborn screening for PKU in 1963; by 1975, all 50 states had implemented mandatory PKU screening. The introduction of tandem mass spectrometry in the 1990s allowed simultaneous screening for PKU and 30--50 other inborn errors of metabolism from a single dried blood spot. The development of sapropterin dihydrochloride (Kuvan), a synthetic form of the natural cofactor tetrahydrobiopterin (BH4) for the treatment of BH4-responsive PKU, was approved by the FDA in 2007. Pegvaliase (Palynziq), a PEGylated recombinant phenylalanine ammonia lyase enzyme for the treatment of adults with PKU not controlled by diet, was approved by the FDA in 2018.

\section*{Why This Test Is Ordered}
Universal newborn screening for PKU: a heel stick blood spot is collected from all newborns at 24--48 hours of life. Diagnostic testing: quantitative plasma amino acid analysis in infants with a positive newborn screen, in older children or adults with unexplained intellectual disability, or in infants born outside of regions with newborn screening. Monitoring of plasma phenylalanine levels in patients with known PKU to assess dietary compliance and adjust dietary therapy. Monitoring during pregnancy in women with PKU (maternal PKU syndrome: if a pregnant woman with PKU has elevated Phe >360 mcmol/L, the fetus is at risk for microcephaly, intellectual disability, congenital heart disease, and intrauterine growth restriction from teratogenic hyperphenylalaninemia).

\section*{Primary vs Secondary Test}
This is a primary newborn screening test for PKU. Universal newborn screening is mandatory in the United States.

\section*{How This Test Is Performed}
Newborn screening: heel stick blood is collected onto filter paper (Guthrie card), air-dried, and sent to the state newborn screening laboratory. Diagnostic testing: venous blood sample for plasma amino acid analysis.

\section*{What Preparation Is Needed}
No special preparation is required for the newborn heel-stick screen. The infant should have received breast milk or formula (protein-containing feeds) for at least 24 hours before collection, because screening before adequate protein intake can produce falsely normal phenylalanine levels. For diagnostic or monitoring plasma amino acid analysis, no fasting is required, although consistent timing relative to feeds is used in monitoring protocols. The sample is collected on a Guthrie card (dried blood spot) or as venous blood and handled per the laboratory's instructions.

\section*{Contraindications and Precautions}
\begin{itemize}
  \item There are no contraindications to newborn screening or venipuncture. Important interpretive cautions:
  \item (1) A screen collected too early, from a premature infant, or before adequate protein feeding can be falsely negative, so a repeat screen is required at 2 weeks of age in these situations.
  \item (2) Elevated phenylalanine is not specific to PKU --- secondary hyperphenylalaninemia can result from tetrahydrobiopterin (BH4) deficiency, liver dysfunction, or protein overload, so confirmatory plasma amino acid analysis and, when indicated, pterin and DHPR testing are essential before starting a restrictive diet.
  \item (3) The same elevated phenylalanine threshold in a woman planning pregnancy requires urgent counseling because maternal hyperphenylalaninemia is teratogenic.
\end{itemize}
\section*{Risks}
Minimal risk from heel stick or venipuncture: slight pain, bruising, or bleeding at the puncture site. Rarely: infection at the site.

\section*{What Happens After the Test}
Newborn screening results are typically available within a few days. A positive screen triggers immediate confirmatory plasma amino acid analysis and referral to a metabolic center; dietary treatment with a phenylalanine-restricted diet and medical formula is initiated without delay when the diagnosis is confirmed. Once diagnosed, patients require lifelong monitoring of blood phenylalanine to keep levels within target ranges, periodic tyrosine measurement, genetic counseling, and for women of reproductive age, strict metabolic control before and during pregnancy to prevent maternal PKU syndrome \cite{blau2010phenylketonuria}.

\section*{Understanding Results}

\subsection*{Normal Newborn Screen}
Plasma phenylalanine <120 mcmol/L (2 mg/dL). Normal phenylalanine/tyrosine ratio. PKU is excluded. No further action required.

\subsection*{Abnormal Newborn Screen (Positive PKU Screen)}
Plasma phenylalanine >120 mcmol/L on newborn screen or confirmatory plasma amino acid analysis. Classical PKU: phenylalanine >1,200 mcmol/L (20 mg/dL) at diagnosis, with dietary phenylalanine tolerance <250--350 mg/day. Mild PKU: phenylalanine 600--1,200 mcmol/L (10--20 mg/dL). Mild hyperphenylalaninemia (non-PKU HPA): phenylalanine 120--600 mcmol/L (2--10 mg/dL) at diagnosis, often managed with less stringent dietary restriction. Treatment: immediate initiation of phenylalanine-restricted diet with phenylalanine-free medical formula. The goal of dietary therapy: maintain blood phenylalanine at 120--360 mcmol/L (2--6 mg/dL) for children <12 years, 120--600 mcmol/L (2--10 mg/dL) for adolescents, and 120--600 mcmol/L for adults. Lifelong dietary adherence is recommended. BH4 responsiveness testing and sapropterin therapy if the patient is BH4-responsive (approximately 20--50\% of PKU patients have some degree of BH4 responsiveness) \cite{vockley2014phenylalanine}.

\subsection*{Maternal PKU (Pregnancy in a Woman with PKU)}
Women with PKU must achieve strict metabolic control (Phe 120--360 mcmol/L) before conception and maintain it throughout pregnancy to prevent maternal PKU syndrome in the fetus. Teratogenic effects: microcephaly (73\%), intellectual disability (92\%), congenital heart disease (12\%), and intrauterine growth restriction (40\%) are strongly correlated with maternal Phe levels during the first trimester \cite{platt2000maternal}.

\section*{Normal Results}
Newborn screening (DBS): Phenylalanine <120 mcmol/L (<2 mg/dL) on a dried blood spot. Confirmatory plasma amino acid analysis: Phenylalanine <120 mcmol/L (2 mg/dL). Phenylalanine/tyrosine ratio <3.0 (molar ratio). Treatment target for PKU: Blood phenylalanine 120--360 mcmol/L (2--6 mg/dL) for children, 120--600 mcmol/L for adolescents and adults.

\section*{Follow-up Tests}
Quantitative plasma amino acid analysis (if the newborn screen is positive). PAH gene sequencing (molecular genetic testing to identify the specific PAH mutation and predict BH4-responsiveness). BH4 loading test (sapropterin responsiveness). Urine pterin profile and DHPR enzyme activity in red blood cells (to distinguish PKU/PAH deficiency from the biopterin synthesis defects [BH4 deficiency] that also cause hyperphenylalaninemia but require a different treatment [BH4, L-dopa, and 5-hydroxytryptophan]). Dietary phenylalanine tolerance testing. Regular monitoring of plasma phenylalanine and tyrosine (every 1--2 weeks in infants, monthly in children, every 1--3 months in adolescents and adults).

\section*{References}
\bibliographystyle{unsrtnat}
\bibliography{references}

\clearpage
\section*{Regulatory Status and Authorization Date}
This chapter describes a test category, not a specific commercial product. FDA, CDSCO, EU IVDR, and MHRA approval or registration dates therefore cannot be assigned without the assay manufacturer, product name, instrument, and intended use. For laboratory-developed tests, an individual agency approval date may not exist. See the Preface for the interpretation of regulatory status.

\clearpage
